\subsection{Persistent components and glass constraints}

\begin{theorem}[Persistent \(H^0\) components only merge]\label{thm:persistent-h0-components-merge}\lean{Infomaterial.Relational.Material.Topology.components_only_merge}\leanok\uses{thm:cover-nerve}
In a graph filtration, path-connected components are monotone: increasing the
graph can merge components but cannot split an existing connection.
\end{theorem}

\begin{theorem}[Kinetically constrained activation is closure]\label{thm:kcm-activation-closure}\lean{Infomaterial.Relational.Material.Glass.activation_is_facilitation_closure}\leanok\uses{thm:causal-settling-transitive-closure}
A kinetically constrained activation process equals the inductive facilitation
closure generated by the seed and the facilitation graph.
\end{theorem}

\begin{theorem}[Finite glass overlap is symmetric]\label{thm:glass-overlap-symmetric}\lean{Infomaterial.Relational.Material.Glass.overlapCount_symm}\leanok\uses{thm:valuation}
The finite replica-overlap count counts the sites where two configurations
agree.  The count is symmetric, and self-overlap fills the whole finite site
list in the companion theorem.
\end{theorem}

\begin{theorem}[Zero mobility is agreement on support]\label{thm:glass-mobility-zero-agreement}\lean{Infomaterial.Relational.Material.Glass.mobility_zero_iff_agree_on_support}\leanok\uses{thm:glass-overlap-symmetric}
The dynamic-heterogeneity count is the finite mobility count.  It vanishes
exactly when two configurations agree on every listed site.
\end{theorem}

\begin{theorem}[Basin descent has no cycle]\label{thm:basin-no-descending-cycle}\lean{Infomaterial.Relational.Material.Glass.BasinGraph.no_descending_cycle}\leanok\uses{thm:no-periodic-orbits}
A finite basin graph with strictly decreasing energy along descent has no
descending cycle. This is the finite landscape H-theorem.
\end{theorem}

\begin{proposition}[Glass substance bundle]\label{thm:glass-substance-bundle}\lean{Infomaterial.Relational.Material.Glass.Substance}\leanok\uses{thm:charge-coloured-sofic, thm:lenticule-hole, thm:homometric-finite-diffraction-mass, thm:kcm-activation-closure, thm:basin-no-descending-cycle, thm:finite-thermodynamic-transition-production}
A glass substance is the bundled material carrier whose layers are
Tiles, Beads, Crystals, Quasicrystals, and the finite glass diagnostics:
configuration overlap, mobility, basin descent, kinetic activation, and the
stationary temperature/Landauer transition.
\end{proposition}

\begin{theorem}[Beads sit between Tiles and Crystals]\label{thm:glass-beads-between-tiles-and-crystals}\lean{Infomaterial.Relational.Material.Glass.Substance.beads_between_tiles_and_crystals}\leanok\uses{thm:glass-substance-bundle}
The glass substance projection from crystals back down to tiles passes through
the bead layer.  This records Beads as a genuine intermediate carrier, not a
prose synonym for either tiles or crystals.
\end{theorem}

\begin{theorem}[A glass substance descends in its basin landscape]\label{thm:glass-substance-descent-lowers-energy}\lean{Infomaterial.Relational.Material.Glass.Substance.descent_lowers_energy, Infomaterial.Relational.Material.Glass.Substance.descent_lowers_rationalEnergy, Infomaterial.Relational.Material.Glass.BasinGraph.rationalEnergy_le_of_descendChain}\leanok\uses{thm:glass-substance-bundle, thm:basin-no-descending-cycle}
The descent witness carried by a glass substance strictly lowers the retained
basin energy, and the same descent lowers the internal rational free-energy
lift of that basin landscape.
\end{theorem}

\begin{theorem}[A glass substance retains stationary temperature]\label{thm:glass-substance-temperature}\lean{Infomaterial.Relational.Material.Glass.Substance.temperature_eq_distribution, Infomaterial.Relational.Material.Glass.Substance.produces_creation_landauer_and_discharge}\leanok\uses{thm:glass-substance-bundle, thm:finite-thermodynamic-transition-production}
The temperature value retained by a glass substance is the stationary
distribution at its distinguished state, and the lifted transition carries
positive creation, balanced flux, positive Landauer cost, and positive
hard-flow discharge.
\end{theorem}

\begin{theorem}[Basin thermalization is glass temperature]\label{thm:glass-basin-thermalization}\lean{Infomaterial.Relational.Material.Glass.Substance.BasinThermalization.temperature_eq_distribution, Infomaterial.Relational.Material.Glass.Substance.BasinThermalization.ledgerAt_eq_eulerChar, Infomaterial.Relational.Material.Glass.Substance.BasinThermalization.produces_landauer_and_discharge}\leanok\uses{thm:glass-substance-temperature, thm:anti-landauer-index-balance, thm:markov-stationary-closure-fixed}
A basin-indexed thermalization puts the stochastic kernel directly on the
metastable basin landscape.  Its temperature is the stationary distribution at
the distinguished basin, each basin carries a Landauer ledger equal to the Euler
characteristic of its support, and the distinguished basin carries balanced flux,
positive Landauer cost, and positive hard-flow discharge.
\end{theorem}

\begin{theorem}[Frozen--hopping threshold]\label{thm:glass-frozen-hopping-threshold}\lean{Infomaterial.Relational.Material.Glass.Substance.BasinThermalization.frozenAt_iff_distribution_le, Infomaterial.Relational.Material.Glass.Substance.BasinThermalization.hoppingAt_iff_le_distribution, Infomaterial.Relational.Material.Glass.Substance.BasinThermalization.threshold_total, Infomaterial.Relational.Material.Glass.Substance.BasinThermalization.frozen_and_hopping_eq_temperature}\leanok\uses{thm:glass-basin-thermalization, thm:real-stochastic-interval-value-category}
The ordered stochastic carrier of a basin thermalization defines frozen and
hopping predicates relative to a threshold.  Stationarity rewrites both
predicates as comparisons with the basin distribution value.  Any total
instance of that order classifies the basin as frozen or hopping, and
antisymmetry identifies a threshold that is both above and below the temperature
with the temperature itself.
\end{theorem}

\begin{theorem}[Basin free-energy bridge]\label{thm:glass-basin-free-energy-bridge}\lean{Infomaterial.Relational.Material.Glass.Substance.BasinFreeEnergyBridge.orbit_Omega_le_init, Infomaterial.Relational.Material.Glass.Substance.BasinFreeEnergyBridge.minimizer_is_equilibrium, Infomaterial.Relational.Material.Glass.Substance.BasinFreeEnergyBridge.temperature_eq_distribution, Infomaterial.Relational.Material.Glass.Substance.BasinFreeEnergyBridge.ledgerAt_eq_eulerChar, Infomaterial.Relational.Material.Glass.Substance.BasinFreeEnergyBridge.produces_temperature_landauer_and_descent}\leanok\uses{thm:glass-basin-thermalization, thm:glass-substance-descent-lowers-energy, thm:positive-temperature-free-energy-flow}
A basin free-energy bridge puts the lower liquid DFT free-energy descent on the
same basin carrier as the stochastic temperature and Landauer ledger.  Along
the basin relaxation orbit, rational free energy stays below its initial value;
minimizers are fixed points; and the same bridge retains stationary temperature,
Euler-characteristic ledger, balanced flux, positive Landauer cost, and positive
hard-flow discharge.
\end{theorem}
